\begin{statement}{E.70}{Proved}{The full Gaussian pro-three source and its compatible norm lifts}
Let $P$ be the free pro-three group on $X,Y$, and form the
actual relative Gaussian triangle
\[
 \begin{gathered}
 \mathfrak G=P\rtimes\langle d\rangle,\qquad
 d^4=1,\quad dXd^{-1}=Y^{-1},\quad dYd^{-1}=X,\\
 x=Xd,\qquad y=Y^{-1}d^2,\qquad z=d.
 \end{gathered}
 \tag{E.70.1}
\]
It is the quotient of the original quartic triangle
obtained by killing its central fourth power and taking
the maximal pro-three quotient of its free kernel.
Let $\psi:\widehat F\to\mathfrak G^6$ evaluate the same
free word in
\[
                (x,y),(y,x),(z,x),(x,z),(y,z),(z,y).
 \tag{E.70.2}
\]
Write $C=P/\gamma_3(P)$ in its integral class-two
coordinates $(v,c)\in\mathbf Z_3^2\times\mathbf Z_3$,
with the law (E.69.1) over $\mathbf Z_3$. Define
\[
 A(v,c)=(-v,c-(v_1+v_2)/2),\qquad
 B(v,c)=(-Sv,c-(v_1+v_2)/2).
 \tag{E.70.3}
\]
Then the full six-coordinate image on the actual source
$N'=N_B'$ is exactly
\[
 \begin{aligned}
 \mathscr H&:=\psi(N')=\psi(\widehat F')\\
 &=\{(g,h)\in P^3\times P^3:
       \overline h=(A(\overline g_3),B(\overline g_2),B(\overline g_1))\}.
 \end{aligned}
 \tag{E.70.4}
\]
Here the bars are the specified maps to $C$. Projection
onto the first three factors is onto $P^3$, with kernel
$\gamma_3(P)^3$ in the last three factors. Moreover
\[
             \mathscr H^{\mathrm{ab}}\simeq\mathbf Z_3^{12}.
 \tag{E.70.5}
\]
The six weight-three coordinates in this abelianization
are retained in the following common-source assertion.

Keep exactly the finite E.66 map $\Psi:N'\to Z$ and
the unchanged E.67 Mordell map $\Theta_M:N'\to\mathcal M$.
The full common image is
\[
       (\Psi,\Theta_M,\psi)(N')
            =Z\times\mathcal M\times\mathscr H.
 \tag{E.70.6}
\]
Every one of the $3456$ finite combinations in E.69
admits a compatible lift through the entire Gaussian
pro-three group and all its characteristic finite
quotients. For each of the four Gaussian actions, the
lift can be chosen to be the actual automorphism
$\operatorname{Inn}(x^j)$, $0\leq j<4$, with both
frames in $P$ and the original six word equations
retained. Thus the finite E.69 frame and all the
specified E.66 and Mordell values can be retained by
one profinite word.

For one chosen full target over an actual E.17 origin
$f_i$, and one such correction $u\in N'$, its entire
word fibre inside $f_iN'$ is
\[
 f_i u\bigl(N'\cap\ker\Psi\cap\ker\Theta_M\cap\ker\psi\bigr).
 \tag{E.70.7}
\]
The Mordell factor here remains its specified order-$54$
quotient. The statement supplies the full Gaussian
tower comparison, with the displayed finite conditions;
it assigns no full-GT lift or additional genus-one
compatibility equation.
\end{statement}
\begin{proof}
We first prove the exact six-coordinate image, then
identify its complete mod-three abelianization and
its actual common quotient with the retained source.

\paragraph{The actual free-kernel maps.}
The presentation (E.70.1) is the exact E.64 quartic
presentation, followed by its indicated quotient.
The source generators still have product $xyz=1$,
and direct free reduction gives
\[
 y^2=z^4=1,\qquad x^4=c:=[X,Y^{-1}],\qquad
 x^2y=X,\qquad z^2y=Y.
 \tag{E.70.8}
\]
In particular $x,y$ generate the whole target
$\mathfrak G$. All six maps (E.70.2) are actual
continuous free-source maps. On $\widehat F'$ they
have deck exponent zero and therefore take values
in $P^6$.

The maps $A,B$ preserve the law of $C$ because
their linear parts have determinant one and their
central corrections are linear. For the three
ordered pairs of evaluations
\[
          (3,4,A),\qquad(2,5,B),\qquad(1,6,B),
 \tag{E.70.9}
\]
there are exact intertwining identities for
conjugation by each of the two free-source
generators. They can be checked without any
finite reduction: conjugation by $(w,d^e)$ on $C$
sends
\[
 (v,c_0)\longmapsto
      (S^ev,c_0+\det(\log(w)_{\mathrm{ab}},S^ev)).
 \tag{E.70.10}
\]
Substitution in (E.70.3) gives all six polynomial
identities over $\mathbf Z[1/2]$. On the normal
generator $[x_{\rm s},y_{\rm s}]$ of $\widehat F'$,
the three comparisons, in the same order as
(E.70.9), are
\[
 \begin{array}{c|c|c}
 &\text{first class-two value}&\text{second class-two value}\\ \hline
 A&(-1,-1,-1/2)&(1,1,1/2)\\
 B&(-1,-1,-1/2)&(1,-1,1/2)\\
 B&(1,1,1/2)&(-1,1,-1/2).
 \end{array}
 \tag{E.70.11}
\]
The six covariance identities propagate these three
identities under all source conjugations, products,
inverses and limits. Since the closed normal closure
of that commutator is $\widehat F'$, they prove the
upper inclusion in (E.70.4) on the whole required
domain. The exact polynomial and reduced-word
checks are saved in \texttt{bc\_gaussian\_pro3\_maps.py}.

The E.67 Gaussian image on $N'$ gives all six
independent coordinates of the first three copies
of $P/\Phi(P)$. Hence its closed image in $P^3$
is onto. Indeed, a proper closed subgroup would
have a proper image in some finite three-group,
and a maximal subgroup of that finite group would
give a missing mod-three abelian character.
The same argument proves surjectivity onto the
last three copies of $P$, and onto the pair of
coordinates $(1,4)$, whose E.67 values are
independently $u$ and $-t$.

\paragraph{The three full kernel factors, with their source membership.}
Here $x_{\rm s},y_{\rm s}$ are the original free-source
generators and $z_{\rm s}=(x_{\rm s}y_{\rm s})^{-1}$.
Literal free reduction in (E.70.1) gives the complete
six-tuples
\[
 \begin{split}
 \psi(x_{\rm s}^{4})&=(c,1,1,c,1,1),\\
 \psi(z_{\rm s}^{4})&=(1,1,1,1,c,dcd^{-1}).
 \end{split}
 \tag{E.70.12}
\]
The elements $x_{\rm s}^4,z_{\rm s}^4$ are used in
$\widehat F$; no assertion that they belong to $N'$
is needed. The subgroup $N$ is normal in the actual
modular source, so $N'$ is normal in $\widehat F$.
Consequently their commutators with any $h\in N'$
do belong to $N'$.

For any $a\in P$, use the last-three-factor
surjection to choose $h\in N'$ with those coordinates
$(1,a,1)$. Then $[z_{\rm s}^4,h]$ has only its fifth
coordinate nontrivial, equal to $[c,a]$. Choosing
$(1,1,a)$ gives only $[dcd^{-1},a]$ in the sixth
coordinate. Finally choose $h\in N'$ with first
coordinate one and fourth coordinate $a$, using
the proved $(1,4)$ surjection. The element
$[x_{\rm s}^4,h]$ has only its fourth coordinate
nontrivial, equal to $[c,a]$.

The closed normal subgroup generated by
$[c,X],[c,Y]$ is exactly $\gamma_3(P)$. Both words
belong to $\gamma_3(P)$. In their quotient $c$ is
central, because $X,Y$ generate $P$. The relation
$c=[X,Y^{-1}]$ then makes $[X,Y]$ central as well;
the quotient is class two, giving the reverse
inclusion. Conjugating $c$ by $d$ gives the same
conclusion since $\gamma_3(P)$ is characteristic.
The kernel of projection of $\psi(N')$ onto the
first three factors is normal under the full
last-three-factor image $P^3$. The three isolated
families above therefore supply every element of
$\gamma_3(P)^3$ in that kernel. The upper inclusion
already bounds the kernel by this group. Together
with the first-three-factor surjection this proves
equality in (E.70.4). It also proves equality with
$\psi(\widehat F')$.

\paragraph{All twelve abelian characters and a finite detector.}
Put $c_0=[X,Y]$. The following continuous
endomorphisms are actual automorphisms of $P$:
\[
 \begin{array}{c|cc}
       &X&Y\\ \hline
 \widetilde A&X^{-1}c_0^{-1/2}&Y^{-1}c_0^{-1/2}\\
 \widetilde B&Yc_0^{-1/2}&X^{-1}c_0^{-1/2}.
 \end{array}
 \tag{E.70.13}
\]
The powers are continuous pro-three powers. Their
Frattini matrices are $-I$ and $-S$, so both maps
are onto. They are isomorphisms: the iterated
Frattini subgroups give cofinal characteristic
finite quotients, on which these surjections are
bijective, and their kernels consequently have
trivial intersection. Their maps on $C$ are
exactly $A,B$.

For the pairs $(1,6),(2,5),(3,4)$, apply the
appropriate lift in (E.70.13) to the first entry.
The second entry times its inverse lies in
$\gamma_3(P)$. This gives actual split descriptions
of the three fibre products as
$\gamma_3(P)\rtimes P$, with conjugation through
the indicated lift. Hence
\[
 \mathscr H^{\mathrm{ab}}
   \simeq\bigl(P^{\mathrm{ab}}\oplus
                    \gamma_3(P)/\gamma_4(P)\bigr)^3
   \simeq(\mathbf Z_3^2\oplus\mathbf Z_3^2)^3.
 \tag{E.70.14}
\]
For completeness the last rank statement follows
from the actual class-three collection coordinates
\[
 X^aY^b c_0^e[c_0,X]^u[c_0,Y]^v,
                     \qquad a,b,e,u,v\in\mathbf Z_3.
 \tag{E.70.15}
\]
Commutator collection gives existence, first for
integer words and then by compactness. Uniqueness
is detected by the corresponding upper-unitriangular
$4$-by-$4$ matrices: the two linear coefficients
detect $a,b$, the weight-two coefficient detects
$e$, and the two independent weight-three
coefficients detect $u,v$. Thus these are ordinary
integral lower-central coordinates, including at
the prime three.

Use the characteristic finite quotient
\[
 P_3=P/\overline{\gamma_4(P)\gamma_2(P)^3P^9}.
 \tag{E.70.16}
\]
In (E.70.15) its coordinate ranges are
$(9,9,3,3,3)$, so $|P_3|=3^7$. The common
class-two quotient has order $3^5$. The image
$S_3$ of $\mathscr H$ in $P_3^6$ is the corresponding
threefold fibre product, of order $3^{27}$.
The splits above reduce to
\[
 S_3\simeq(P_3\times\mathbf F_3^2)^3,
 \qquad
 \mathscr H/\Phi(\mathscr H)
            \xrightarrow{\ \sim\ }S_3/\Phi(S_3)
            \simeq\mathbf F_3^{12}.
 \tag{E.70.17}
\]
One can see the last map directly: in each pair
retain the first entry's two abelian coordinates
modulo three and the two weight-three coordinates
of the second entry times the inverse of the
specified lift of the first. Formula (E.70.15)
shows that these are precisely the four reductions
of its integral abelianization coordinates.

The finite certificate implements all these maps.
Write a strictly upper-unitriangular matrix as
$(a_{12},a_{23},a_{34},a_{13},a_{24},a_{14})$.
The four pairs of generator images in $U_4(\mathbf F_3)$ are
\[
 \begin{array}{c|c|c}
 &X&Y\\ \hline
 1&(1,1,1,0,0,0)&0\\
 2&0&(1,1,1,0,0,0)\\
 3&(1,1,0,0,0,0)&(0,0,1,0,0,0)\\
 4&(1,0,0,0,0,0)&(0,1,1,0,0,0).
 \end{array}
 \tag{E.70.18}
\]
Their $2187$ collected values are all distinct.
Four central characters, evaluated on
$X^3,Y^3,[c_0,X],[c_0,Y]$, have rank four.
Each matrix acts faithfully
on its three affine coordinates, and the product
action has degree $108$. Induction through the
actual order-four deck automorphism gives degree
$432$, and all six literal evaluations give degree
$2592$. The program checks the deck automorphism
on every element and both generating products,
with the right-action convention
$p(gh)=p(h)[p(g)]$. A separate character certificate
decodes the actual six matrices and checks all
twelve values of (E.70.17), including every saved
word DAG; their rank is twelve.

\paragraph{The exact common image with the retained finite source.}
Adjoin those six actions and the six unchanged
Mordell actions to the actual $5223$-point E.66
source. The source has degree $7869$. All $97$
literal $N$ words are compared with the E.67
export. The complete Schreier calculation first
computes the image of $N$, and then its derived
subgroup, which is the image of $N'$. Only after
this step does it retain the five coordinate-zero
distribution blocks. The resulting degree is
$7495$. The independent full Schreier replays give
\[
 \begin{array}{c|c}
 \text{actual retained projection}&\text{order}\\ \hline
 S_3&3^{27}\\
 \mathcal M\times S_3&3^{30}\\
 Z\text{ and }S_3&|Z|\,3^{27}\\
 Z\text{ and }\mathcal M\text{ and }S_3&|Z|\,3^{30}.
 \end{array}
 \tag{E.70.19}
\]
Each equality is for the actual common-domain map;
it follows from the complete source certificates,
not a comparison of independent projection orders.

Since (E.70.17) is the full Frattini quotient,
the closed image of $N'$ in $\mathcal M\times\mathscr H$
is the whole pro-three product. A common quotient
of that product with finite $Z$ is a finite
three-group. Any nontrivial such quotient has a
nontrivial cyclic quotient of order three. Its
character factors through the specific map
\[
 \mathcal M\times\mathscr H
       \longrightarrow\mathcal M\times S_3
       \longrightarrow
          (\mathcal M\times S_3)/\Phi(\mathcal M\times S_3).
 \tag{E.70.20}
\]
This would give a common cyclic quotient in the
last actual product in (E.70.19), which is
impossible. Goursat's exact quotient construction
therefore proves (E.70.6).

\paragraph{Compatible full Gaussian frames for every retained finite value.}
For $j=0,1,2,3$, let $\alpha_j=\operatorname{Inn}(x^j)$.
It fixes $x$ and its primitive fourth power.
The elements $y$ and $\alpha_j(y)$ have order two
and the same image in the deck group. Coprime
complement conjugacy gives a frame $h_y\in P$.
The element $h_z=x^jd^{-j}$ is already a frame for
$\alpha_j(z)$. Their possible class-two central
coordinates may be prescribed independently.
Here is the exact lifting argument: for an element
fixed by $y$ or $d$ in $C$, lift it to $P$. Its
conjugation sends the order-two or order-four
complement to another complement differing within
$\gamma_3(P)$. Conjugacy within that pro-three
kernel corrects the lift. This proves the
surjections of the corresponding fixed centralizers.
For a finite three-group kernel, the conjugacy
assertion follows by induction through a central
invariant elementary-abelian subgroup and averaging
the resulting cocycle, since two and four are units.
For the pro-three kernel, the finite conjugator sets
are nested nonempty compact sets. This proves the
claimed conjugacy and centralizer maps in the full
groups.

The fixed centralizer in $C$ consists of its central
coordinate in each case: $-I-I$ and $S-I$ are
invertible over $\mathbf Z_3$. We can therefore
choose the frames to have precisely
\[
 \begin{array}{c|c|c}
 j&\overline h_y&\overline h_z\\ \hline
 0&(0,0,q)&(0,0,0)\\
 1&(1/2,1/2,q)&(1,0,1/4)\\
 2&(1,0,q)&(1,-1,0)\\
 3&(1/2,-1/2,q)&(0,-1,-1/4),
 \end{array}
 \qquad q\in\mathbf Z_3.
 \tag{E.70.21}
\]
The abelian coordinates follow by conjugating the
literal elements (E.70.1). Substitution of these
coordinates in the six values (E.69.12) satisfies
all three integral graph identities (E.70.4).
The four polynomial checks hold for the unrestricted
parameter $q$ and are included in the literal-map
certificate. Thus the entire six-tuple lies in
$\mathscr H$, while its full group action is the
chosen actual $\alpha_j$. Reducing (E.70.21) modulo
three gives exactly the twelve E.69 framed targets.

For an actual origin $f_i$, its six Gaussian values
lie in $\mathscr H$ by (E.70.4). The correction from
them to any of the chosen full norm tuples is
therefore an element of $\mathscr H$. Choose its
simultaneous $N'$ preimage in (E.70.6), retaining
the prescribed E.66 correction and the unique
finite Mordell frame. This proves the compatible
whole-tower lift of every E.69 finite combination.
The exact kernel of that common homomorphism gives
the full word fibre (E.70.7).
\end{proof}
